\documentclass{article}
\usepackage{amsmath,amssymb}
\usepackage{xurl}
\usepackage[hidelinks]{hyperref}
\setlength{\emergencystretch}{2em}
% Shared commands for notes in docs/paper-gaps/.

\DeclareUnicodeCharacter{2080}{\ifmmode{}_{0}\else\textsubscript{0}\fi}
\DeclareUnicodeCharacter{2081}{\ifmmode{}_{1}\else\textsubscript{1}\fi}
\DeclareUnicodeCharacter{2082}{\ifmmode{}_{2}\else\textsubscript{2}\fi}
\DeclareUnicodeCharacter{2099}{\ifmmode{}_{n}\else\textsubscript{n}\fi}

\newcommand{\Lean}{\textsc{Lean}}
\ExplSyntaxOn
\NewDocumentCommand{\leanid}{m}{
  \ifmmode
    \textsf{\detokenize{#1}}
  \else
    \group_begin:
    \urlstyle{sf}
    \tl_set:Nn \l_tmpa_tl {#1}
    \tl_replace_all:Nnn \l_tmpa_tl {\_} {_}
    \exp_args:NV \path \l_tmpa_tl
    \group_end:
  \fi
}
\ExplSyntaxOff
\providecommand{\tr}{\operatorname{tr}}
\newcommand{\tnleanIssueBase}{https://github.com/LionSR/TNLean/issues/}
\newcommand{\tnleanPrBase}{https://github.com/LionSR/TNLean/pull/}
\newcommand{\tnleanDocBase}{https://sirui-lu.com/TNLean/docs/}
\newcommand{\ghissue}[1]{\href{\tnleanIssueBase#1}{GitHub issue \##1}}
\newcommand{\ghpr}[1]{\href{\tnleanPrBase#1}{PR \##1}}
\newcommand{\doclink}[1]{\href{\tnleanDocBase#1.html}{\path{#1}}}

% Dirac notation and the identity operator, as in blueprint/src/macros/common.tex.
% \providecommand keeps note-local definitions authoritative.
\usepackage{dsfont}
\providecommand{\ket}[1]{|#1\rangle}
\providecommand{\bra}[1]{\langle #1|}
\providecommand{\braket}[2]{\langle #1 | #2 \rangle}
\providecommand{\ketbra}[2]{|#1\rangle\!\langle #2|}
\providecommand{\Id}{\mathds{1}}
\providecommand{\one}{\mathds{1}}
\providecommand{\C}{\mathbb{C}}
\providecommand{\MN}[1]{M_{#1}(\C)}
\providecommand{\MD}{M_D(\C)}

% Verdict marker: \gapnote{<kind>}{<status>}. Kind is the policy
% classification (clarification, local-correction, scope-restriction,
% unfaithful, false-source, open-gap); status is open, wip (elimination
% underway), resolved, or historical. Machine-read by the site index;
% prints a verdict line.
\newcommand{\gapnote}[2]{\par\noindent\textbf{Verdict:} #1 (#2).\par}

\title{The Honeycomb Cluster State: Sign of the Parent Hamiltonian and Injectivity}
\author{}
\date{2026-11-18}
\begin{document}
\maketitle
\gapnote{local-correction}{open}

\section{The source claim}
Section~IV of \cite{Cirac2021Matrix} introduces the cluster state as follows:
\begin{quote}
    ``The canonical example for an injective PEPS with global symmetries is
    the cluster state on the honeycomb lattice, which is the unique ground
    state of the commuting parent Hamiltonian $\sum_{ijkl}X_iZ_jZ_kZ_l$ of
    qubits where $j,k,l$ are the nearest neigbours of site $i$. This state
    has a very simple PEPS representation: [figure] with [vertex tensor]
    $=\ket{0000}+\ket{1111}$, [bond matrix]
    $=\frac{1}{\sqrt{2}}\left(\begin{smallmatrix} 1 & 1\\ 1 &
                -1\end{smallmatrix}\right)$.
    It exhibits the global symmetry $Y^{\otimes N}$, as this is just the
    product of all (commuting) terms of the Hamiltonian. To use the
    fundamental theorem of PEPS, we first block two sites of this PEPS such
    as to make it uniform [\ldots]''
\end{quote}
\footnote{Source traceability: \path{Papers/2011.12127/TN-Review-main.tex},
    lines 1445--1452.}

\section{What is formalized}
For a finite simple graph with vertex set $V$ and edge set $E$, put a copy
tensor $\ket{0\ldots0}+\ket{1\ldots1}$ at every vertex and the matrix
$H=\frac1{\sqrt2}\left(\begin{smallmatrix}1&1\\1&-1\end{smallmatrix}\right)$
on every bond. Contracting the network gives the coefficients
\begin{equation*}
    \psi(\sigma)=2^{-|E|/2}\prod_{\{u,v\}\in E}(-1)^{\sigma_u\sigma_v},
\end{equation*}
which is $2^{(|V|-|E|)/2}$ times the controlled-$Z$ circuit applied to
$\ket{+}^{\otimes V}$. In the formal tensor the bond matrix is absorbed into
the leg at the smaller endpoint of each bond, which leaves the contraction
unchanged. The stabilizers $K_i=X_i\prod_{j\sim i}Z_j$ commute and square
to the identity, the operator $-\sum_iK_i$ is bounded below by $-|V|$, and
the eigenspace for $-|V|$ is the line spanned by $\psi$. These statements
are specialized to a brick-wall honeycomb on the torus, where every site
has three nearest neighbours.\footnote{\leanid{TNLean.PEPS.stateCoeff_honeycombClusterPEPS},
    \leanid{TNLean.PEPS.honeycombStabilizer_comm},
    \leanid{TNLean.PEPS.honeycombHamiltonian_eq_neg_card_smul_iff},
    \leanid{TNLean.PEPS.honeycombGraph_degree}, in
    \doclink{TNLean/PEPS/Examples/HoneycombCluster}.}

\section{The sign of the parent Hamiltonian}
The cluster state $\ket{C}$ satisfies $K_i\ket{C}=\ket{C}$ for every $i$.
Since the $K_i$ commute and square to the identity, the spectrum of
$\sum_iK_i$ is contained in $\{-|V|,\ldots,|V|\}$, and its lowest eigenvalue
$-|V|$ is attained exactly on the common $-1$ eigenspace. That space is
spanned by $\prod_iZ_i\ket{C}$: the operator $Z_i$ anticommutes with $K_i$
and commutes with $K_j$ for $j\neq i$. So the printed Hamiltonian
$+\sum X_iZ_jZ_kZ_l$ has the unique ground state $\prod_iZ_i\ket{C}$, which
is not the cluster state, although the two are related by a product of
on-site unitaries. The other statements of the passage are insensitive to
the sign: the product of all terms is $\pm\prod_i X_iZ_i^3$, a multiple of
$Y^{\otimes N}$ either way.

\textbf{Local fix.} The parent Hamiltonian is taken to be
$-\sum_iX_iZ_jZ_kZ_l$, the usual stabilizer convention, under which the
source's claim holds. In Lean this is \leanid{TNLean.PEPS.honeycombHamiltonian}.

\section{Injectivity}
The vertex tensor $\ket{0000}+\ket{1111}$ maps the virtual space of a
vertex of degree three, of dimension eight, to the qubit, and its image is
two-dimensional. A vertex tensor with at least two bonds is therefore not
injective, with or without the bond matrices, and the honeycomb PEPS is not
injective site by site.\footnote{\leanid{TNLean.PEPS.not_isVertexInjective_honeycombClusterPEPS}.}
Blocking the two sites of a unit cell, as the source does to make the
tensor uniform, does not help: the blocked tensor has four virtual legs of
dimension two and a physical space of dimension four.

The claim holds after blocking over hexagonal plaquettes. A plaquette has
six sites, each with exactly one bond leaving the plaquette. With the
boundary indices $\beta$ fixed, the blocked tensor sends $\beta$ to
\begin{equation*}
    \prod_{\text{inner bonds}}CZ\;\bigotimes_{u}M_u\ket{\beta_u},
\end{equation*}
up to normalization, where each $M_u$ is the identity or $H$ according to
the side on which the bond matrix of the leaving bond is placed. Both maps
are invertible, so the blocked tensor is injective. The honeycomb vertex
set is partitioned by every third plaquette, on a torus whose dimensions
are compatible with this pattern, and the blocked network is a triangular
lattice of injective tensors. This argument is not formalized, and no
injectivity of the honeycomb PEPS is asserted.

\section{The torus}
The brick wall on a torus of width $w$ and height $h$ joins each site to
its two horizontal neighbours and to one vertical neighbour, above when the
coordinate sum is even and below otherwise. For even $w$ and $h$ it is the
honeycomb lattice. For even $h$ and $w\ge3$ every site has exactly three
distinct neighbours; for $w=2$ the two horizontal neighbours coincide. The
general graph statements hold for every finite simple graph, so they apply
to the brick wall on every torus of width and height at least two.

\section{Elimination plan}
\begin{enumerate}
    \item Formalize injectivity of the plaquette-blocked tensor and the
          partition of the honeycomb torus into plaquettes, then state the
          source's injectivity under that blocking.
    \item The sign is a convention and needs no further work.
\end{enumerate}

\bibliographystyle{alpha}
\bibliography{references}
\end{document}
